\section{Implementation}
\label{sec:implementation}

\subsection{Technology Stack}

\software{} is implemented as an npm monorepo under the \texttt{app/} directory.
Table~\ref{tab:techstack} summarizes the technology choices and their
rationale.

\begin{table}[htbp]
  \centering
  \caption{Technology stack summary.}
  \label{tab:techstack}
  \begin{tabularx}{\textwidth}{L{3cm} L{4cm} L{7cm}}
    \toprule
    \textbf{Layer} & \textbf{Choice} & \textbf{Rationale} \\
    \midrule
    Frontend framework & React 18~\cite{reactjs} & Component ecosystem, TypeScript support,
    large community \\
    Bundler & Vite~\cite{vite} & Fast HMR, native TypeScript/JSX support \\
    Editor widget & CodeMirror 6~\cite{codemirror} & Extensible, \LaTeX{} syntax highlighting,
    autocomplete \\
    Math rendering & KaTeX~\cite{katex} & Fast browser-side math, no server dependency \\
    PDF preview & Native browser PDF embed & In-place compiled PDF viewing without an
    additional rendering runtime \\
    HTTP server & Fastify~\cite{fastify} & Schema validation, serialization, plugin system \\
    AI streaming & SSE~\cite{sse_spec} & Bidirectional-free streaming, HTTP compatibility \\
    Testing & Vitest~\cite{vitest} & Fast, Vite-compatible test runner \\
    Build system & npm workspaces and concurrently & Workspace orchestration and
    coordinated frontend/backend development scripts \\
    Terminal & tmux~\cite{tmux} & Session persistence, pane management \\
    \bottomrule
  \end{tabularx}
\end{table}

\subsection{Project Management and File Operations}

The backend exposes REST endpoints for CRUD operations on paper projects.
Projects are identified by directory name under \texttt{papers/}, with metadata
stored in \texttt{project.json}. The project listing endpoint includes
auto-detection logic: when it encounters a directory containing recognizable
paper files (\texttt{.tex}, \texttt{.bib}, \texttt{.pdf}, \texttt{.sty},
\texttt{.cls}, \texttt{.md}) without \texttt{project.json}, it generates
compatible metadata automatically. This enables drag-and-drop import of
existing \LaTeX{} projects without manual configuration.

File operations (shown below) are implemented through a service
layer that validates paths against the project root to prevent directory
traversal attacks. The file tree API returns the complete directory structure
as a JSON tree, which the frontend renders as an interactive component with
context menus supporting copy, cut, paste, rename, delete, and drag-and-drop
file moves. Path validation uses path normalization and prefix checking:

\begin{lstlisting}[language=JavaScript, caption={Path security validation pattern}]
function resolveProjectPath(projectRoot, userPath) {
  const normalized = path.normalize(userPath);
  const resolved = path.resolve(projectRoot, normalized);
  // Prevent directory traversal outside project root
  if (!resolved.startsWith(projectRoot)) {
    throw new Error('Path traversal detected');
  }
  return resolved;
}
\end{lstlisting}

Auto-detection of paper projects also scans for multiple file types:
\texttt{.pdf} files indicate compiled output, \texttt{.sty} and \texttt{.cls}
indicate style templates, \texttt{.bib} indicates bibliography databases, and
\texttt{.md} indicates Markdown notes or supplementary material. The system
also supports projects containing only Markdown files, enabling use cases
beyond strict \LaTeX{} workflows.

\subsection{Multi-Mode Editor and Preview}

The editor component supports three view modes built on CodeMirror 6:

\begin{itemize}
  \item \textbf{Source mode:} Standard code editor with \LaTeX{} syntax
  highlighting, BibTeX citation autocomplete (typing \texttt{@} triggers a
  search query against CrossRef for real-time DOI lookup), and real-time
  document model.

  \item \textbf{Split mode:} Source editor alongside a rendered preview panel.
  The preview resolves project-relative image references through the project
  blob API, so user-created folders such as \texttt{fig/}, \texttt{figures/},
  \texttt{images/}, or \texttt{img/} display correctly inside the editor
  preview without browser-page-relative URL resolution.

  \item \textbf{Rendered mode:} An editable preview surface inspired by
  Obsidian's live preview~\cite{obsidian}. Valid Markdown/\LaTeX{} blocks are compiled into
  rendered text, math, lists, and images. Edits to rendered text blocks are
  written back to the source file, while syntax that cannot be rendered stays
  visible as editable source fallback.
\end{itemize}

The preview layer implements a custom \LaTeX{} subset renderer covering common
constructs: sections, abstracts, lists, quotes, verbatim, theorem/proof
environments, tables, math environments (including \texttt{equation},
\texttt{align}, \texttt{gather}), code listings, citations, and cross-references.
This design compromise---supporting a curated subset rather than full \LaTeX{}
rendering---was driven by the observation that most manuscript editing revolves
around these common constructs, and attempting full \LaTeX{} rendering in the
browser would be prohibitively complex (competing with server-side \TeX{}
engines).

\subsection{AI Integration and Permission Modes}

The AI service abstracts over provider-specific APIs (OpenAI chat completions,
Anthropic Messages API) behind a unified interface with configurable model,
base URL, and system prompt. Configuration is loaded from the repository-root
\texttt{.env} file, and the frontend settings panel writes changes back through
\texttt{/api/config}. API keys are never exposed to the browser---the config
endpoint returns only masked status (e.g., \texttt{sk-****}).

\textbf{Streaming.} AI responses are delivered via Server-Sent Events (SSE),
with event types for tokens, tool calls, tool results, completion, and errors.
The frontend renders tokens incrementally for a typewriter effect. Automatic
fallback to non-streaming mode ensures compatibility when SSE connections
fail, such as behind certain proxy configurations.

\textbf{Context injection.} Before each AI call, the backend assembles a
context package based on the conversation's scope setting:
\begin{itemize}
  \item \textbf{Chapter scope:} Reads the full content of the selected project
  file (e.g., \texttt{sec/3.design.tex}) and injects it as context.
  \item \textbf{Global scope:} Reads the first 400 characters of each
  \texttt{.tex} file to provide overview context without token overflow.
  \item \textbf{Reference scope:} Injects the complete \texttt{references.bib}
  content for citation-aware operations.
\end{itemize}
All non-free scopes additionally inject \texttt{references.bib} content.

\textbf{Permission mode enforcement.} Each conversation is assigned one of
three modes at creation time:
\begin{itemize}
  \item \textbf{Chat mode:} The system prompt includes instructions to discuss
  only; no tool definitions for file writing or code execution are provided to
  the model.
  \item \textbf{Agent mode:} The model receives tool definitions for
  \texttt{edit\_file}, which returns a unified diff. The diff is displayed
  inline with color-coded additions and deletions, and the user must
  explicitly Accept or Reject each change.
  \item \textbf{Tools mode:} The model receives full tool access, including
  shell execution, but tools are scoped to the \texttt{code/} project
  directory. All tool invocations and their results are recorded in the
  conversation history.
\end{itemize}

\subsection{Pipeline Engine}

Pipeline 2.0 defines a composable stage system with five typed executors,
managed by a pipeline engine that handles state persistence, stage transitions,
and error recovery. Pipeline state is serialized to
\texttt{\textasciitilde/.paper-writer/pipelines/} for durability across server
restarts.

\begin{itemize}
  \item \textbf{AI stages:} Execute LLM-powered operations using skill plugins.
  Skills are modular prompts that define a system message, output format, and
  evaluation criteria. Over 40 skills are available, covering proofreading,
  section outlining, abstract generation, argument analysis, and style
  adjustment.
  \item \textbf{Compute stages:} Execute shell commands with configurable
  timeouts (default 30 seconds). Standard output and error streams are captured
  and displayed in the pipeline log.
  \item \textbf{Human stages:} Pause pipeline execution and present a
  checkpoint interface with options to approve, reject, provide feedback,
  edit inline, or skip. Human checkpoints proved to be a frequently used pipeline
  feature in practice.
  \item \textbf{Citation stages:} Execute batch citation verification with
  rate-limited API calls (respecting crossref.org rate limits with exponential
  backoff and jitter).
  \item \textbf{Compile stages:} Execute \LaTeX{} compilation via the
  configured engine (pdflatex, xelatex, lualatex, latexmk, or tectonic), parse
  the \texttt{.log} file for errors and warnings, and return structured error
  reports.
\end{itemize}

\subsection{Citation Verification}

The citation service implements a multi-strategy verification approach
(Figure~\ref{fig:citation-strategy}):

\begin{enumerate}
  \item \textbf{DOI-based verification:} Look up DOIs against
  CrossRef~\cite{crossrefapi} and Semantic Scholar~\cite{semanticscholarapi}
  to confirm existence, retrieve correct title, authors, and publication venue.
  \item \textbf{Title-based fuzzy matching:} For entries without DOIs, search
  by title against CrossRef and OpenAlex~\cite{openalex} using approximate
  string matching to handle minor title variations.
  \item \textbf{Local cross-checking:} Parse \texttt{.tex} files for
  \texttt{\textbackslash cite\{\}} commands and compare against BibTeX keys in
  \texttt{.bib}, flagging missing or undefined references.
  \item \textbf{Hallucination screening:} Cross-reference all citation keys
  against known scholarly identifiers to flag potentially hallucinated
  references.
\end{enumerate}

\begin{figure}[htbp]
  \centering
  \includegraphics[width=\textwidth]{fig-citation-strategy.pdf}
  \caption{Citation verification strategy flowchart. The verification
  process begins by parsing the \texttt{.bib} file into individual entries.
  For each entry, the system first checks whether a DOI is present. \emph{Entries
  with DOIs} are verified directly against the CrossRef REST API and the
  Semantic Scholar API by resolving the DOI and comparing the returned metadata
  (title, authors, year) against the bibliography entry. \emph{Entries without
  DOIs} use title-based fuzzy matching against CrossRef and OpenAlex to identify
  candidate records. Confidence levels are assigned based on cross-API agreement:
  \textbf{high} (2+ APIs confirm the entry with matching metadata), \textbf{medium}
  (1 API confirms), \textbf{low} (title match only with metadata discrepancies),
  or \textbf{none} (unverifiable, flagged for manual review). Arrows indicate
  the decision flow; dashed arrows indicate fallback paths when primary
  verification fails.}
  \label{fig:citation-strategy}
\end{figure}

Confidence levels are assigned based on API agreement: \textbf{high} (2+ APIs
confirm), \textbf{medium} (1 API confirms), \textbf{low} (title match only),
or \textbf{none} (unverifiable). For large bibliographies (100+ entries), batch
processing with rate limiting and parallelization completes within a few
minutes.

\subsection{Model Context Protocol Integration}

Building on the MCP design introduced in Section~\ref{sec:design},
\software{} implements the specification as a JSON-RPC 2.0 server with SSE
transport. Seven tools are registered, shown in Table~\ref{tab:mcp-tools}:

\begin{table}[htbp]
  \centering
  \caption{MCP tool registry: the seven tools exposed by the \software{} MCP
  server. Each tool is accessible via JSON-RPC 2.0 over SSE transport from any
  MCP-compatible client (Claude Desktop, Cursor, GitHub Copilot).}
  \label{tab:mcp-tools}
  \begin{tabularx}{\textwidth}{L{4cm} L{10cm}}
    \toprule
    \textbf{Tool Name} & \textbf{Description} \\
    \midrule
    \texttt{search\_papers} & Search papers by title, keyword, or DOI \\
    \texttt{verify\_citation} & Verify a citation against scholarly databases \\
    \texttt{cross\_check\_citations} & Cross-check all citations in a project \\
    \texttt{compile\_latex} & Compile \LaTeX{} project to PDF \\
    \texttt{read\_project\_file} & Read a file from the project tree \\
    \texttt{polish\_text} & Polish academic text via AI \\
    \texttt{review\_paper} & Run AI review on a paper \\
    \bottomrule
  \end{tabularx}
\end{table}

The service discovery endpoint returns tool schemas in JSON Schema format,
enabling auto-configuration of MCP clients. The implementation requires
approximately 200 lines of server code, demonstrating the low integration cost
of the MCP standard.

Two concrete integration scenarios were used during development. First, Claude
Desktop was configured with the Paper Agent MCP endpoint and used to invoke
\texttt{cross\_check\_citations} on a manuscript directory. The client received
the citation-check schema from service discovery, requested the project path,
and returned missing-key and bibliography-status diagnostics without requiring
the user to copy \texttt{.tex} and \texttt{.bib} contents into a chat window.
This converted a manual workflow---open source files, search for
\texttt{\textbackslash cite\{\}}, compare against \texttt{references.bib}, and
summarize missing entries---into a single tool call backed by the same verifier
used in the Paper Agent UI.

Second, Cursor was configured against the same MCP endpoint and used to invoke
\texttt{compile\_latex} from an editor conversation. The tool executed the
configured \LaTeX{} engine, parsed the compilation log, and returned structured
errors to the client.

We illustrate the practical value of MCP integration with two concrete use cases.

\paragraph{Use case 1: Citation verification via Claude Desktop.}
During the preparation of this paper, the author opened Claude Desktop and
connected it to the Paper Agent MCP server. A natural-language prompt---``Check
all citations in my bibliography against CrossRef and report any that cannot be
verified''---invoked the \texttt{verify\_citations} tool. The tool parsed the
project's \texttt{.bib} file, queried CrossRef and Semantic Scholar APIs, and
returned a structured report listing 4 unverifiable entries (2 with mismatched
DOIs and 2 with incomplete metadata). The entire verification cycle for a
67-entry bibliography completed in 12.3 seconds. Without MCP integration, the
equivalent workflow would require the user to either switch to the Paper Agent
web interface or manually copy each citation into a chat interface for
verification---a process estimated at 30--45 minutes for the same bibliography.
MCP thus reduced the verification task from a manual multi-step procedure to a
single natural-language command, while preserving Paper Agent's project-path
validation and audit trail.

\paragraph{Use case 2: \LaTeX{} compilation and error diagnosis via Cursor.}
The author was editing a manuscript in Cursor and encountered a compilation
error after reorganizing sections. Rather than switching to a terminal to run
\texttt{pdflatex} and inspect the log, the author typed: ``Compile my paper
and tell me what's wrong.'' Cursor invoked the \texttt{compile\_latex} MCP
tool, which executed the configured engine (pdflatex), captured the log, parsed
the errors, and returned a structured summary: two undefined references caused
by a missing \texttt{\textbackslash label\{\}} and one missing package
(\texttt{booktabs}). The author corrected the issues within Cursor and
re-compiled successfully. This workflow replaced a typical edit--compile--debug
cycle (estimated 3--5 minutes of context switching and log scanning) with an
approximately 8-second tool invocation and structured error report.

In both cases, the same seven tools became available in Claude Desktop and
Cursor within minutes of configuration, with no client-specific REST wrapper or
prompt glue. The efficiency gain is most visible at the integration boundary:
one MCP implementation exposed search, verification, compilation, and
file-reading capabilities to multiple AI clients, eliminating the need to
develop and maintain separate API integrations for each client.

\subsection{Collaboration Infrastructure}

\software{} includes preliminary collaboration infrastructure built on
WebSocket-based operational merge and a token-based authorization system.
Cryptographically signed tokens with configurable time-to-live (default 24
hours) allow sharing read-write or read-only project access.

The merge mechanism works as follows. The collaboration layer maintains a
server-side document snapshot for each shared file, stored as a plain-text
string in memory. When a client sends a text-change operation over WebSocket
(encoded as a JSON message containing the file path, change offset, deleted
length, and inserted text), the server (1) validates the token and project
path, (2) applies the operation to the current snapshot at the specified
offset, (3) broadcasts the updated content to all subscribed clients, and (4)
writes the result to disk through a debounced flush mechanism (800 ms window)
that batches multiple rapid edits into a single filesystem write, reducing
storage I/O and preventing partial-write corruption. Token validation uses
HMAC-SHA256 signatures to prevent unauthorized access or token tampering.
This design supports basic use cases such as co-author review sessions and
advisor feedback rounds without requiring participants to share filesystem
access.

The mechanism is intentionally simpler than full operational transform or CRDT
synchronization~\cite{crdt2011}. The current collaboration infrastructure does
\emph{not} provide the following capabilities:
\begin{itemize}
  \item Live cursor presence (showing other users' cursor positions).
  \item Colored selections indicating which user is editing which region.
  \item Semantic conflict resolution for simultaneous edits to the same text
  range (e.g., OT-based intention preservation).
  \item Offline edit reconciliation (edits made while disconnected cannot be
  merged automatically upon reconnection).
  \item Undo/redo that is aware of multi-user edit history.
\end{itemize}
\noindent Conflicts are handled conservatively by last-writer-wins snapshot
replacement: when two clients edit the same region concurrently, the later
operation overwrites the earlier one. This approach is acceptable for
short review sessions where collaborators edit different sections, but it is
unsuitable for heavy multi-author simultaneous drafting. Full real-time
collaborative editing similar to Overleaf's multi-user experience remains
future work (see Section~\ref{sec:limitations}).

\subsection{Terminal and Code Execution}

The integrated terminal is backed by tmux sessions mapped per-project. Each
project directory maps to a stable tmux session name, so reopening the terminal
or refreshing the page reattaches to the same shell state. File upload to the
project directory is supported through the transfer sub-agent, which handles
format detection and directory-aware placement.

Code execution within the workspace is managed through a code executor service
with configurable sandbox support. The service supports command execution with
timeout management, output capture, and termination signaling. Synthetic
terminal interaction via pty.js enables running interactive tools that require
pseudo-terminal allocation.

\subsection{Writing Pattern Analysis}

The system includes an \textbf{experimental} writing pattern analysis panel
that computes surface-level statistics---perplexity estimates, sentence-length
variance, vocabulary richness, and repetition patterns---from the current
manuscript. These metrics are compared against baseline distributions derived
from published academic prose to flag passages that may trigger automated
AI-detection heuristics.

\textbf{Important limitations and disclaimer.} This feature is \emph{not}
designed for, and must \emph{not} be used as, a tool for evading
academic-integrity checks or AI-detection systems. AI detection tools have
significant false-positive and false-negative rates~\cite{liang2024llmwriting},
and their outputs are unreliable indicators of whether text was human-written
or machine-generated. The writing-pattern panel is framed in the interface as a
\emph{diagnostic aid for self-review}---helping authors identify passages that
may read as atypically uniform or mechanical---rather than as an authorship
classifier or a detector-evasion tool. The interface displays explicit warnings
that: (1) the panel's metrics are heuristic and discipline-dependent; (2) no
score can prove that text is human-written or machine-written; (3) the feature
must not be used to circumvent institutional plagiarism or AI-use policies; and
(4) authors bear sole responsibility for the authenticity and integrity of
their submitted work. Its limitations are further discussed in
Section~\ref{sec:limitations}.

\subsection{Testing and Quality Assurance}

The repository's automated testing infrastructure comprises 32 Vitest test
files plus a Playwright end-to-end specification covering:
\begin{itemize}
  \item \textbf{Backend unit tests:} Project routes, service behavior,
  configuration parsing and privacy, conversation store, compile service,
  citation verification, pipeline stage types and executors, terminal
  management, and file manager operations.
  \item \textbf{Frontend unit tests:} Project tree creation logic and UI
  interaction, rendered editor mode, skill engine integration, layout and
  resize behavior.
  \item \textbf{Integration tests:} API endpoint integration, conversation
  mode switching, pipeline compilation workflow, citation verification
  pipeline.
  \item \textbf{E2E tests:} Browser-based project lifecycle (Playwright).
\end{itemize}

Tests are executed via Vitest, with configuration supporting both JavaScript
and TypeScript files. The test suite was designed to grow incrementally with
each feature addition, following a red-green-refactor cycle. The frontend
production build serves as an additional integration smoke test.
